\section{引言}

月球反射通信（Earth-Moon-Earth, EME）以月球表面作为被动反射体，可在超远距离上建立链路。Evans 的雷达月球观测给出了月球回波的早期测量[1]。在单载波 BPSK 体制下，长回波表现为数十至上百符号的记忆：本文使用的 Level B 配置最大时延为 116 符号，标称相干时间约 1200 s。深衰落与长记忆叠加，使接收端必须在低信噪比下同时完成信道估计、相位跟踪与符号判决。

传统方案以 Pilot 驱动的线性最小均方误差有限冲激响应（LMMSE-FIR）均衡器与判决反馈均衡器（DFE-RLS）为代表。它们在短记忆、慢变信道中有效，但在 116 符号记忆下，前缀 Pilot 可提供的独立方程少于信道抽头数，最小二乘问题欠定；判决反馈又把低信噪比下的错误判决送回滤波器，形成误差传播。本文主矩阵中，0 dB 下 CFO+DD-Phase DFE-RLS 的 Data BER 为 44.05\%，LMMSE-FIR 为 44.01\%，已接近随机猜测。

神经均衡器提供了另一条路径。ViterbiNet 用神经网络替代 Viterbi 算法中的分支度量[2]；深度展开把迭代估计器的每一步映射为网络层[3]；Pilot 条件化的接收机与元学习进一步减少了推理时所需的 Pilot[4,5]。这些工作说明，离线仿真数据可以训练出强均衡器。本文主矩阵中，离线神经均衡器在 0/5/10/15 dB 的 Data BER 为 22.19\%/5.96\%/1.06\%/0.216\%，相对 CFO+DD-Phase DFE-RLS 降低 49.6\%–95.1\%。三个独立 seed 在每个 SNR 上均给出同向结论。

仿真训练与部署的差异不仅来自噪声实现，还来自 CIR 抽头、相位与 CFO 的随机实现。一个在训练分布内很强的网络，遇到分布外信道时可能同时出现均衡残差与置信度失准；如果只依赖离线权重，接收机没有机制纠正这种失配。在线适配因此是长记忆信道下值得优先考虑的环节。问题在于适配什么：网络权重的维数高、在线缺少标签、更新容易失稳；物理状态的维数低、由 Pilot 与接收波形直接观测，更适合作为在线自由度。

离线训练依赖仿真信道与完整标签，因此上述失配在部署时难以避免。本文因此不把网络权重作为在线自由度。一个更可控的自由度是物理状态：CIR 的维数远低于网络参数，且可由 Pilot 与判决反馈持续观测。于是问题变为：在冻结神经权重的条件下，仅递推 CIR 后验，能否获得可归因、可重复的 Data BER 增益？该增益是否随会话内帧数累积？

本文围绕这两个问题做三件事。第一，给出一个冻结神经块均衡器加贝叶斯 CIR 后验的接收机：离线训练 Pilot 条件化的展开网络；在线只更新 116 维 CIR 的均值与协方差，Data 位置使用模型软判决作为伪观测，并以符号方差构造行噪声，使不可靠符号自动降权。第二，在运行前冻结配置、超参数、checkpoint 与判据，用 10 个独立 session × 128 帧做一次性确认，统计单位为 session 而非帧或 bit。第三，分别报告主效应、逐 SNR 安全护栏与增长假设：H1 与 H2 通过，H3 在 5/10/15 dB 通过、0 dB 不通过。本文的在线自由度仅限物理状态，所有增益均与同源 Frozen 网络配对比较。
